% ----------------------------------------------------------
% Introdução 
% 
% ----------------------------------------------------------

\chapter[Introdução]{Introdução}

A metabolômica é uma das áreas mais recentes das ciências ômicas e tem como objetivo identificar e quantificar os metabólitos presentes em células, tecidos, fluidos biológicos e organismos. Os metabólitos são pequenas moléculas, geralmente com massa inferior a 1500 Da, e correspondem aos produtos finais das reações bioquímicas que acontecem no organismo. Enquanto a genômica analisa o DNA, a transcriptômica analisa o RNA e a proteômica analisa as proteínas, a metabolômica observa o resultado final desse fluxo de informação: o que de fato está sendo produzido pelas células em um determinado momento \cite{qiu2023small}. Por essa proximidade com o funcionamento real do organismo, alterações no perfil metabólico têm sido associadas a doenças, à resposta a tratamentos e a fatores nutricionais, ambientais e farmacológicos, o que tornou a metabolômica uma abordagem promissora para a descoberta de biomarcadores e para o avanço da medicina de precisão.
 
Entre as plataformas analíticas disponíveis, a cromatografia líquida acoplada à espectrometria de massas (LC-MS) é hoje uma das mais utilizadas em estudos de metabolômica não direcionada, abordagem que busca detectar o maior número possível de sinais na amostra, sem definir antecipadamente quais metabólitos analisar \cite{schrimpe2016untargeted}. A técnica funciona em duas etapas: primeiro, a cromatografia líquida separa os compostos da amostra com base em propriedades físico-químicas, como polaridade e carga; depois, o espectrômetro de massas detecta esses compostos e fornece informações como massa acurada, padrão isotópico e fragmentação molecular. Cada sinal detectado é descrito por três parâmetros (tempo de retenção, razão massa/carga (\textit{m/z}) e intensidade) e é chamado de \textit{feature} metabolômica, unidade básica das matrizes de dados que serão analisadas nas etapas seguintes \cite{chen2022guide}.
 
Apesar do grande potencial da LC-MS, os dados gerados por essa técnica não são simples de analisar. Os arquivos brutos contêm milhares de sinais, que precisam ser organizados em uma matriz estruturada por meio de etapas como alinhamento cromatográfico, detecção de picos e agrupamento de adutos, fragmentos e isótopos. Para essa etapa inicial, é comum o uso de softwares comerciais como o Progenesis QI, da Waters/Nonlinear Dynamics. Mesmo depois desse processamento primário, porém, a matriz resultante ainda pode apresentar variações técnicas entre amostras, sinais pouco reprodutíveis, redundâncias analíticas e falsos positivos causados por ruído do equipamento ou contaminação ambiental \cite{schiffman2019filtering}.
 
Por causa dessas inconsistências, antes de qualquer análise estatística é necessário realizar etapas adicionais de pré-processamento, como a avaliação da qualidade dos dados a partir de amostras de controle de qualidade (QCs), a filtragem de \textit{features} com alto desvio padrão relativo (RSD), o tratamento dos valores ausentes, a normalização entre amostras e a transformação e o escalonamento das intensidades. Para cada uma dessas etapas existem diferentes métodos disponíveis, e a escolha entre eles não é trivial. Cada decisão pode influenciar diretamente quais metabólitos aparecem como mais relevantes no resultado final, fazendo com que uma mesma matriz, processada de maneiras diferentes, possa levar a conclusões distintas \cite{vandenberg2006centering}. A complexidade aumenta ainda mais na análise estatística subsequente, dado que os dados metabolômicos têm alta dimensionalidade (com número de variáveis muito maior do que o número de amostras) e exigem métodos multivariados específicos, como PCA, PLS-DA e Random Forest \cite{worley2013multivariate}.
 
Outro problema comum na metabolômica é a fragmentação do fluxo de trabalho: cada etapa da análise costuma ser feita em um software diferente, como Progenesis QI, XCMS, MetaboAnalyst, SIMCA, R ou Python, o que aumenta o risco de erros, dificulta a padronização dos procedimentos e faz com que o resultado dependa muito da experiência individual do pesquisador \cite{weber2017computational}. Diante desse cenário, o uso de inteligência artificial (IA) e de aprendizado de máquina vem se mostrando uma alternativa interessante para apoiar o pesquisador nas decisões metodológicas que envolvem muitas variáveis técnicas e biológicas ao mesmo tempo. Quando essas técnicas são reunidas em um único ambiente de trabalho, é possível reduzir a fragmentação atual e tornar as decisões analíticas mais transparentes, padronizadas e reprodutíveis \cite{sen2021deep}.
 
É nesse contexto que se insere o presente trabalho, que propõe o desenvolvimento de uma plataforma para assistir no refino quimiométrico de matrizes de dados de LC-MS. A ferramenta proposta recebe como entrada as matrizes semiprocessadas geradas em softwares como o Progenesis QI e conduz, de forma assistida, etapas como filtragem por RSD, tratamento de valores ausentes, normalização, transformação, escalonamento e análise estatística.

\section{Objetivos}

\subsection{Objetivo Geral}

Desenvolver uma plataforma para o refino quimiométrico assistido de matrizes semiprocessadas de dados de metabolômica por cromatografia líquida acoplada à espectrometria de massas (LC-MS).

\subsection{Objetivos Específicos}

\subsubsection{Sistematizar as etapas de pré-processamento de dados metabolômicos não direcionados}
 
Identificar e organizar as principais etapas do refino quimiométrico (filtragem por desvio padrão relativo (RSD), tratamento de valores ausentes, normalização, transformação e escalonamento), que são determinantes para a interpretação biológica dos resultados em metabolômica \cite{vandenberg2006centering}.



\subsubsection{Estruturar um fluxo de trabalho assistido a partir das matrizes exportadas pelo Progenesis QI}
 
Desenvolver um pipeline integrado que receba as matrizes semiprocessadas geradas pelo Progenesis QI e conduza, de forma assistida, as etapas seguintes do refino quimiométrico, evitando a necessidade de alternar entre várias ferramentas distintas durante a análise.
 

\subsubsection{Incorporar algoritmos como sistema de apoio à decisão metodológica}
 
Integrar algoritmos determinísticos e de aprendizado de máquina capazes de orientar, a partir das características dos dados de entrada, a escolha dos métodos mais adequados de pré-processamento, normalização e análise estatística \cite{galal2022applications}.


\subsubsection{Implementar funcionalidades de análise estatística univariada e multivariada com critérios de validação}
 
Disponibilizar testes estatísticos univariados, com correção para múltiplas comparações via FDR, e métodos multivariados como PCA, PLS-DA e Random Forest, acompanhados de estratégias de validação cruzada e testes de permutação para reduzir o risco de sobreajuste \cite{worley2013multivariate}.

\subsubsection{Avaliar o desempenho da plataforma em dados reais e sintéticos}
 
Validar a ferramenta aplicando-a aos dados de metabolômica não direcionada do laboratório parceiro e a matrizes sintéticas com defeito conhecido, comparando os resultados com os do pipeline em R usado anteriormente pelo laboratório e avaliando a consistência das decisões metodológicas sugeridas pelo sistema.

\section{Justificativa}
 
A análise de dados metabolômicos não direcionados por LC-MS envolve diversos desafios técnicos, estatísticos e de interpretação. Mesmo depois de passar por um processamento inicial em softwares como o Progenesis QI, as matrizes resultantes ainda contêm variações técnicas, valores ausentes, sinais pouco reprodutíveis, redundâncias analíticas e possíveis falsos positivos. Por isso, são necessárias etapas adicionais de refino quimiométrico antes que qualquer análise estatística possa ser feita de forma confiável \cite{mahieu2017systems}. Essas etapas, no entanto, costumam ser realizadas em ferramentas separadas e frequentemente exigem conhecimentos avançados em estatística e programação. Isso fragmenta o fluxo de trabalho, prejudica a reprodutibilidade dos estudos e faz com que o resultado dependa fortemente da experiência individual do pesquisador \cite{weber2017computational}.
 
Além disso, a escolha entre os diferentes métodos de filtragem, imputação, normalização, transformação, escalonamento e análise estatística não é trivial. Como cada decisão pode levar a resultados diferentes a partir de uma mesma matriz, é fundamental contar com critérios sistemáticos, justificados e reprodutíveis \cite{diguida2016non}. Nesse cenário, ferramentas baseadas em algoritmos de auxílio à decisão, capazes de processar grandes volumes de dados, reconhecer padrões e auxiliar em decisões metodológicas, surgem como uma alternativa interessante para tornar o fluxo de análise metabolômica mais acessível, transparente e padronizado \cite{sen2021deep}.
 
Assim, desenvolver uma plataforma integrada para o refino quimiométrico assistido de dados de LC-MS pode contribuir tanto para o avanço metodológico da área quanto para tornar essas análises mais acessíveis, beneficiando pesquisadores com diferentes níveis de experiência em estatística e computação.

\section{Delimitação do estudo}

Este trabalho se limita ao desenvolvimento de uma plataforma computacional baseada em algoritmos de auxílio à decisão, voltada ao refino quimiométrico assistido de matrizes de dados gerados por cromatografia líquida acoplada à espectrometria de massas (LC-MS), no contexto de estudos de metabolômica não direcionada. Outras técnicas analíticas, como ressonância magnética nuclear (NMR), cromatografia gasosa acoplada à espectrometria de massas (GC-MS) e eletroforese capilar acoplada à espectrometria de massas (CE-MS), bem como estudos de metabolômica direcionada, não fazem parte do escopo deste trabalho.

A ferramenta atua em um estágio intermediário do fluxo de análise metabolômica, posicionada entre o processamento primário, feito pelo Progenesis QI, e as etapas finais de interpretação biológica. Sua entrada corresponde às matrizes semiprocessadas exportadas pelo Progenesis QI, normalmente em formato CSV, que contêm a caracterização das amostras (com a definição dos grupos experimentais) e uma planilha estruturada de \textit{features} identificadas por tempo de retenção, razão massa/carga (\textit{m/z}) e intensidade. Estão fora do escopo, portanto, todas as etapas anteriores a essa exportação, como o desenho experimental, a coleta e o preparo das amostras, a aquisição dos dados pelo espectrômetro de massas, a conversão dos arquivos brutos (\textit{raw}) e o alinhamento cromatográfico inicial.

Funcionalmente, a plataforma contempla a avaliação da qualidade analítica da matriz, a filtragem por desvio padrão relativo (RSD), o tratamento de valores ausentes, a normalização, a transformação, o escalonamento e as análises estatísticas univariadas e multivariadas. O foco está nas operações de caráter geral do refino de metabolômica não direcionada por LC-MS, que não dependem de um estudo específico. A validação, porém, é feita com os dados de um único estudo do laboratório parceiro, e com matrizes sintéticas derivadas deles; a generalização para outros estudos, laboratórios e equipamentos não foi avaliada e fica como trabalho futuro. Customizações para projetos individuais e a aplicação clínica direta em pacientes também ficam fora do escopo deste trabalho.